% !TeX root = ../../Book2.tex
% 纲要标题：### C.3 环的 Brauer 群
% 纲要位置：工作记录/计划纲要.md，第 1820 行。

% 纲要要点（供目录核对）：
% * 环的 Brauer 群及其与域情形的区别
% #### C.3.1 等价关系与群运算
% #### C.3.2 分裂类、底环变换与域情形

\section{环的 Brauer 群}
\label{b2:sec:C03}

域的 Brauer 群把增加矩阵大小视为无关变化。一般交换环上的有限投射模可能没有统一的自由基，所以相应的可忽略因子应从矩阵代数扩展为 \(\operatorname{End}_R(P)\)，其中 \(P\) 忠实、有限生成且投射。上一节已经证明，这些因子确实属于 Azumaya 代数。

\subsection{等价关系与群运算}
\label{b2:subsec:C03:01}

\begin{proposition}\label{b2:prop:azumaya-tensor-opposite}
设 \(R\) 为非零交换含幺环。若 \(A,B\) 为 Azumaya \(R\) 代数，则 \(A^{\mathrm{op}}\) 与 \(A\otimes_RB\) 也为 Azumaya \(R\) 代数。
\end{proposition}
\begin{proof}
反代数不改变底模，张量积的忠实有限投射性由\cref{b2:lem:azumaya-projective-base-change}保证。对任意极大理想 \(\mathfrak m\)，这些代数的纤维分别为 \(A(\mathfrak m)^{\mathrm{op}}\) 与 \(A(\mathfrak m)\otimes_{\kappa(\mathfrak m)}B(\mathfrak m)\)。原纤维中心单；反代数保留单性与中心，张量积的中心单性由\cref{b2:thm:csa-tensor-product}保证。因此\cref{b2:thm:azumaya-fiber-criterion}适用于两个新代数。
\end{proof}

\begin{definition}\label{b2:def:ring-brauer-equivalence}
设 \(R\) 为非零交换含幺环、\(A,B\) 为 Azumaya \(R\) 代数。若存在忠实有限生成投射 \(R\) 模 \(P,Q\)，以及 \(R\) 代数同构
\begin{equation}\label{b2:eq:ring-brauer-equivalence}
A\otimes_R\operatorname{End}_R(P)
\cong B\otimes_R\operatorname{End}_R(Q),
\end{equation}
则称 \(A,B\) 在 \(R\) 上 Brauer 等价。
\end{definition}

此处同构必须保持指定的 \(R\) 标量，不能只要求抽象环同构。\chapref{b2:ch:18}已经定义域上的 Brauer 等价；这里使用同一名称，并将在下一个小节说明两种定义在域上重合。

\begin{theorem}\label{b2:thm:ring-brauer-group}
设 \(R\) 为非零交换含幺环。在 Azumaya \(R\) 代数上，\cref{b2:def:ring-brauer-equivalence}所定义的 Brauer 等价是等价关系。记 \([A]\) 为代数 \(A\) 的等价类。这些等价类在运算
\[
[A]+[B]=[A\otimes_RB]
\]
下构成交换群，零元为 \([R]\)，\([A]\) 的逆元为 \([A^{\mathrm{op}}]\)。这个群称为交换环 \(R\) 的 Brauer 群，记为 \(\operatorname{Br}(R)\)。
\end{theorem}
\symidx{\(\operatorname{Br}(R)\)：交换环的 Brauer 群}
\begin{proof}
取 \(P=Q=R\) 得自反性，交换同构两端得对称性。设
\[
A\otimes_R\operatorname{End}_R(P)\cong B\otimes_R\operatorname{End}_R(Q),
\qquad
B\otimes_R\operatorname{End}_R(U)\cong C\otimes_R\operatorname{End}_R(V).
\]
将第一式张量 \(\operatorname{End}_R(U)\)，第二式张量 \(\operatorname{End}_R(Q)\)，再交换中间因子，由式\eqref{b2:eq:azumaya-end-tensor}得到
\[
A\otimes_R\operatorname{End}_R(P\otimes_RU)
\cong C\otimes_R\operatorname{End}_R(V\otimes_RQ).
\]
\cref{b2:lem:azumaya-projective-base-change}保证这里的两个张量模仍忠实有限生成且投射，所以关系具有传递性。

若 \(A\) 与 \(A'\)、\(B\) 与 \(B'\) 分别通过式\eqref{b2:eq:ring-brauer-equivalence}等价，将两式张量并交换因子，再用同一个自同态张量公式，就得到 \(A\otimes_RB\) 与 \(A'\otimes_RB'\) 等价。因此运算不依赖代表。结合性、交换性和单位元分别由张量积的结合、交换和单位同构给出。

最后，\(A\) 自身是忠实有限生成投射模，定义中的同构给出
\[
A\otimes_RA^{\mathrm{op}}\cong\operatorname{End}_R(A).
\]
而 \(\operatorname{End}_R(A)\) 与 \(R\) 等价：在等价关系中分别选 \(R\) 和 \(A\) 作为附加模即可。所以 \([A]+[A^{\mathrm{op}}]=[R]\)，每个元素都有逆元。所有代数均以某个有限自由模的直和因子为底模，它们的同构类型构成集合，因而所述等价类确实构成一个群。
\end{proof}

特别地，\(\bigl[\operatorname{End}_R(P)\bigr]=0\)。这一步说明为何必须消去自同态代数，而不只消去固定大小的矩阵环。在不连通的基环上，\(P\) 的秩甚至可能在不同部分取不同的值。

\subsection{分裂类、底环变换与域情形}
\label{b2:subsec:C03:02}

设 \(R\) 为非零交换含幺环、\(A\) 为 Azumaya \(R\) 代数。若存在忠实有限生成投射 \(R\) 模 \(P\)，使 \(A\cong\operatorname{End}_R(P)\) 为 \(R\) 代数同构，则称 \(A\) 在 \(R\) 上分裂。\cref{b2:thm:ring-brauer-group}已经说明分裂代数代表零类。要证明反向，需从等价关系给出的附加自同态因子中恢复 \(A\)；下面的中心化子计算提供了这个恢复方法。

\begin{lemma}[自同态因子的中心化子]\label{b2:lem:azumaya-end-centralizer}
设 \(R\) 为非零交换含幺环、\(A\) 为结合含幺 \(R\) 代数，其结构同态的像位于中心。设 \(P\) 为忠实有限生成投射 \(R\) 模，令 \(E=\operatorname{End}_R(P)\)。在 \(A\otimes_RE\) 中记
\[
C=\{\,x\in A\otimes_RE\mid x(1\otimes e)=(1\otimes e)x\text{ 对全部 }e\in E\text{ 成立}\,\}.
\]
则 \(a\mapsto a\otimes\operatorname{id}_P\) 给出 \(R\) 代数同构 \(A\cong C\)。
\end{lemma}
\begin{proof}
上述映射的像包含在 \(C\) 中。由\cref{b2:lem:azumaya-projective-base-change}，\(E\) 有限生成。选其有限个 \(R\) 模生成元 \(e_1,\ldots,e_v\)，则 \(C\) 是映射
\[
A\otimes_RE\longrightarrow (A\otimes_RE)^v,
\qquad x\longmapsto\bigl(x(1\otimes e_i)-(1\otimes e_i)x\bigr)_{i=1}^v
\]
的核。局部化保持正合性，而且 \(e_i/1\) 生成每个局部化后的自同态代数，故取这个中心化子与局部化相容。

对任意极大理想 \(\mathfrak m\)，\cref{b2:lem:azumaya-local-toolkit,b2:lem:azumaya-projective-base-change}给出 \(P_{\mathfrak m}\cong R_{\mathfrak m}^n\)、\(n\ge1\)，以及
\[
(A\otimes_RE)_{\mathfrak m}\cong M_n(A_{\mathfrak m}).
\]
在这个识别下，与 \(1\otimes E\) 交换等价于与全部矩阵单位交换。与对角矩阵单位交换先使非对角条目为零，与其余矩阵单位交换再使对角条目相等。因此中心化子恰为 \(aI_n\)、\(a\in A_{\mathfrak m}\)，原映射的局部化就是 \(A_{\mathfrak m}\cong C_{\mathfrak m}\)。对原映射的核与余核应用\cref{b2:lem:azumaya-local-toolkit}的局部消失判据，得 \(A\cong C\)。映射公式保留乘法、单位和 \(R\) 标量，故它是 \(R\) 代数同构。
\end{proof}

\begin{proposition}\label{b2:prop:azumaya-zero-split}
设 \(R\) 为非零交换含幺环、\(A\) 为 Azumaya \(R\) 代数。\(A\) 代表 \(\operatorname{Br}(R)\) 中的零类，当且仅当存在忠实有限生成投射 \(R\) 模 \(T\)，使 \(A\cong\operatorname{End}_R(T)\) 为 \(R\) 代数同构。
\end{proposition}
\begin{proof}
若 \(A\cong\operatorname{End}_R(T)\)，\cref{b2:thm:ring-brauer-group}给出 \([A]=0\)。反过来，设 \([A]=0\)，按等价关系选忠实有限生成投射 \(R\) 模 \(P,Q\)，使
\[
A\otimes_RE\cong\operatorname{End}_R(Q),
\qquad E=\operatorname{End}_R(P).
\]
经这个同构，\(Q\) 成为左 \(E\) 模，且 \(E\) 中的 \(R\) 标量仍按原来的 \(R\) 模结构作用于 \(Q\)。

令 \(P^*=\operatorname{Hom}_R(P,R)\)，赋予 \(P\) 双模结构 \({}_EP_R\)，其作用为 \(ep=e(p)\)、\(pr=rp\)；赋予 \(P^*\) 双模结构 \({}_RP^*_E\)，其作用为 \((r\lambda)(p)=r\lambda(p)\)、\((\lambda e)(p)=\lambda(e(p))\)。对偶保留有限自由模的直和因子，所以 \(P^*\) 有限生成且投射，双对偶求值映射 \(P\rightarrow P^{**}\) 为同构。取式\eqref{b2:eq:azumaya-trace-unit}中的 \(p_i,\lambda_i\)；若 \(rP^*=0\)，则 \(r=\sum_i(r\lambda_i)(p_i)=0\)。因此 \(P^*\) 忠实，由\cref{b2:lem:azumaya-projective-base-change}还是生成元。取对偶给出 \(\operatorname{End}_R(P^*)\cong E^{\mathrm{op}}\)，将\cref{b2:cor:morita-inverse-bimodules}用于左 \(R\) 模 \(P^*\)，再用双对偶同构，得到双模同构
\[
{}_RP^*\otimes_EP_R\cong{}_RR_R,
\qquad
{}_EP\otimes_RP^*_E\cong{}_EE_E.
\]
因此 \(P\otimes_R-\) 与 \(P^*\otimes_E-\) 分别从左 \(R\) 模到左 \(E\) 模、从左 \(E\) 模到左 \(R\) 模，且由张量结合律互为拟逆。令 \(T=P^*\otimes_EQ\)，便有 \(Q\cong P\otimes_RT\)，以及由等价的全忠实性得到的 \(R\) 代数同构
\[
\operatorname{End}_E(Q)\cong\operatorname{End}_R(T).
\]
由\cref{b2:lem:azumaya-end-centralizer}，\(A\otimes_RE\) 中 \(1\otimes E\) 的中心化子为 \(A\)；在所给同构下，它恰好对应 \(\operatorname{End}_E(Q)\)。所以 \(A\cong\operatorname{End}_R(T)\)。

还须核对 \(T\) 的模性质。对 \(p\in P\)、\(\lambda\in P^*\)，记 \(\theta_{p,\lambda}\in E\) 为算子 \(x\mapsto\lambda(x)p\)。式\eqref{b2:eq:azumaya-trace-unit}给出右 \(E\) 模同态
\[
E^t\longrightarrow P^*,\qquad (e_i)\longmapsto\sum_i\lambda_i e_i,
\]
其截面为 \(\lambda\mapsto(\theta_{p_i,\lambda})_i\)，因为 \(\sum_i\lambda_i\theta_{p_i,\lambda}=\lambda\)。这些映射也保持 \(R\) 作用，取 \(-\otimes_EQ\) 后使 \(T\) 成为 \(Q^t\) 的 \(R\) 模直和因子。因此 \(T\) 有限生成且投射。若 \(rT=0\)，由 \(Q\cong P\otimes_RT\) 得 \(rQ=0\)，而 \(Q\) 忠实，故 \(r=0\)。这证明 \(T\) 忠实。
\end{proof}

零类因而恰好由分裂代数代表。不过，这些代数之间仍可能不同构；Brauer 等价消去了自同态因子，不能据此恢复原代数的大小。改变底环则会进一步改变哪些类能够分裂。

\begin{proposition}\label{b2:prop:ring-brauer-base-change}
设 \(R,S\) 为非零交换含幺环。环同态 \(f:R\rightarrow S\) 诱导群同态
\[
f^*:\operatorname{Br}(R)\longrightarrow\operatorname{Br}(S),
\qquad [A]\longmapsto[S\otimes_RA].
\]
这些同态保持恒等映射，并按底环变换次序保持复合。若 \(R=k\) 为域，则 \(\operatorname{Br}(k)\) 与\cref{b2:thm:brauer-group}构造的群相同。
\end{proposition}
\begin{proof}
\cref{b2:thm:azumaya-fiber-criterion}保证底环变换后的代数仍为 Azumaya 代数。对式\eqref{b2:eq:ring-brauer-equivalence}取 \(S\otimes_R-\)，式\eqref{b2:eq:azumaya-hom-base-change}把附加因子化为
\[
S\otimes_R\operatorname{End}_R(P)
\cong\operatorname{End}_S(S\otimes_RP),
\]
且 \(S\otimes_RP\) 仍忠实有限生成且投射。故等价关系保持。底环变换与张量积相容，并把 \(R\) 送到 \(S\)，所以所得映射为群同态；两次底环变换的结合识别给出复合相容性。

在域上，忠实有限生成投射模恰为非零有限维向量空间，其自同态代数恰为 \(M_n(k)\)、\(n\ge1\)。又由\cref{b2:prop:azumaya-field-matrix}，代数对象恰为中心单代数。因此等价关系变为 \(M_n(A)\cong M_m(B)\)，群运算也同为张量积，两个构造完全一致。
\end{proof}

例如，取 \(R=\mathbb R[t]\)，由 \(\mathbb R\rightarrow R\) 扩张得到 \(A=R\otimes_{\mathbb R}\mathbb H\)。这是 Azumaya \(R\) 代数，却不代表零类。若 \([A]=0\)，再按 \(t\mapsto0\) 底环变换到 \(\mathbb R\)，便有 \([\mathbb H]=0\)。\cref{b2:prop:brauer-division-representative}说明域上每个 Brauer 类有唯一的中心除代数代表，\(\mathbb H\ne\mathbb R\)，所以这不可能。另一方面，扩张到 \(\mathbb C[t]\) 后 \(A\) 变成 \(M_2\bigl(\mathbb C[t]\bigr)\)，其类随之变为零。

\ProseParagraphBreak
一般交换环上的 Brauer 群仍以代数及其模为对象，但域上“每类选一个除代数”的描述已不适用。例如 \(R=k\times k\) 的任何忠实代数都包含非零的中心幂等元 \((1,0)\)、\((0,1)\)，二者乘积为零，因而不可能是除环。环上的分类问题还须顾及各处的纤维如何在同一个有限投射模上相容。以自同态代数为零类、以张量积为加法的表述，也可参阅 Weibel 的《The K-book: An Introduction to Algebraic K-theory》\cite[Chapter II, Example 5.4.3]{Weibel2013KBook}；该处从对称幺半范畴的群化解释这组关系。
